%% kapa: exam=1 round=25 year=2014 period=2 form=A subject=appraisal_law subject_ko=부동산관계법규 seq=34 number=34
%% source: Q-Net 감정평가사 기출문제 (cst00309) artlSeq=5101207 file=extracted/1008/5101207/2014년 제25회 감정평가사 1차시험 기출문제/2교시(A형).hwp
%% conversion: kordoc 4.24.51 -> tools/gen1.py (auto); review=proofread-sonnet; flags=none
\begin{question}{34}
측량ㆍ수로조사 및 지적에 관한 법령상 부동산종합공부에 관한 설명으로 옳지 않은 것은?
\begin{choices}
  \item 지적소관청은 부동산종합공부를 영구히 보존하여야 한다.
  \item 지적공부의 내용 중 토지의 소유자에 관한 사항은 부동산종합공부의 등록사항이다.
  \item 토지이용계획확인서의 내용 중 토지의 이용 및 규제에 관한 사항은 부동산종합공부의 등록사항이다.
  \item 부동산종합증명서를 발급받으려는 자는 지적소관청 이외에 읍ㆍ면ㆍ동의 장에게도 신청할 수 있다.
  \item 부동산종합공부의 등록사항에 잘못이 있는 경우에는 지적소관청의 직권정정만 허용된다.
\end{choices}
\end{question}
\answer{5}
